\section{A rational certificate for the comparison bound}
\label{cert:section}

We now construct a smooth family of six-component laws and a polynomial
that lies below the attainable value on its boundary.  The polynomial has
slightly more curvature than the lower-test inequality permits at an interior
contact point.  This forces it below the attainable value throughout the
family.  Its value at the center is strictly greater than $\log 3$.

There are two numerical difficulties.  First, the chart is only approximately
harmonic, so we must retain the supporting-covector term in
Theorem~\ref{thm:ap}.  Second, a finite list of point evaluations does not
prove a differential or boundary inequality on a continuous region.  We
address the first by moving the boundary slightly into the positive simplex,
and the second by coefficient majorants, polynomial interval bounds, and
interpolation of actual attainable constructions.  All decimal constants
in this section that define polynomials or arithmetic thresholds denote exact
rational numbers.

\begin{theorem}[The certificate bound]\label{thm:certificate}
Let $\beta=1129/500$.  The construction class defining $F_\beta$ satisfies
\[
 F_\beta(\alpha_*)\ \geq\ 1.098619383270680538
     \ >\ \log 3,
\]
where $\alpha_*=(p,p,p,1/3-p,1/3-p,1/3-p)$ and
\[
 p=2^{-500}\left\lfloor 2^{500}\frac{18623}{100000}
                                      +\frac12\right\rfloor.
\]
\end{theorem}

The proof occupies this section.  The finite checks are specified together
with analytic bounds that turn their integer conclusions into uniform
inequalities.  Section~\ref{cert:reproducibility} gives the implementation
and reproduction interface.  In particular, the computations used below are
finite certificate verifications; they do not presume convergence of a
numerical optimizer or of an infinite harmonic series.

\subsection{The chart and its exact coefficients}
\label{cert:chart}

Put $\zeta=\exp(2\pi i/3)$ and $R=31/25$.  For a real polynomial $P$ on the
complex plane define, with subscripts modulo three,
\begin{equation}\label{cert:law}
 r_i(z)=\frac{1+\Re(\zeta^i z)}3,\qquad
 p_i(z)=P(\zeta^i z),\qquad
 E_i(z)=r_i(z)+p_i(z)-p_{i+1}(z)-p_{i+2}(z).
\end{equation}
Their six-coordinate sum is one.  In the six-state support, their mean
satisfies $\mu_i=1-r_i$.  Thus
\[
 z=3r_0-1+i\sqrt3(r_2-r_1)
\]
is an invertible affine coordinate on the mean plane.  Whenever all six
coordinates are positive, \eqref{cert:law} is a chart to which
Theorem~\ref{thm:ap} applies.

\begin{lemma}[Complex covariance normalization]\label{ap:covariance}
For $\alpha=(p_0,p_1,p_2,E_0,E_1,E_2)$, set
\[
 r_i=1-\mu_i,\qquad
 z=x+iy=3r_0-1+i\sqrt3(r_2-r_1),\qquad
 D_i=r_i(1-r_i)+2p_i.
\]
Then $D_i=\operatorname{Var}(g_i)$.  In the $(x,y)$ coordinates
the covariance contraction has matrix
\begin{subequations}\label{cert:operator}
\begin{equation}\label{ap:real-covariance}
 C=\begin{pmatrix}9D_0&3\sqrt3(D_1-D_2)\\
 3\sqrt3(D_1-D_2)&-3D_0+6D_1+6D_2\end{pmatrix}.
\end{equation}
Writing
\[
 d=\frac{D_0+D_1+D_2}{3},\qquad
 f=\frac{D_0+\zeta D_1+\zeta^2D_2}{3},
\]
the operator is
\begin{equation}\label{ap:complex-operator}
 \mathcal L u
 =36\bigl(d\,u_{z\bar z}+f\,u_{zz}
                         +\bar f\,u_{\bar z\bar z}\bigr).
\end{equation}
\end{subequations}
\end{lemma}
\begin{proof}
Since $g_i$ takes values zero, one, and two, with probability $p_i$
of the last, $\E g_i^2=\mu_i+2p_i$.  Hence
$\operatorname{Var}(g_i)=\mu_i(1-\mu_i)+2p_i=D_i$.
For the centered grade $\xi=g-\mu$, the sum of the three coordinates
is zero, so
\[
 \E\xi_i\xi_j=\frac{D_k-D_i-D_j}{2}
 \quad (\{i,j,k\}=\{0,1,2\}).
\]
The induced coordinate increment is
$\Delta z=-3\xi_0+i\sqrt3(\xi_1-\xi_2)$.
Substitution gives \eqref{ap:real-covariance}, and consequently
\[
 \E|\Delta z|^2=18d,\qquad \E(\Delta z)^2=36f.
\]
For Wirtinger derivatives
$\partial_z=(\partial_x-i\partial_y)/2$ and
$\partial_{\bar z}=(\partial_x+i\partial_y)/2$, a directional
derivative is $\Delta z\,\partial_z+\Delta\bar z\,\partial_{\bar z}$.
Squaring it and taking expectations proves \eqref{ap:complex-operator}.
In particular the coefficient 36 belongs to the full contraction used
in Theorem~\ref{thm:ap}; it is not the generator divided by two.
\end{proof}

Here is a finite rational definition of $P$.  For integers $N\geq8$ and
$b\geq1$, put $Q=2^b$ and $\operatorname{rnd}(t)=\lfloor t+1/2\rfloor$.
Let
\[
 (Z_0,\ldots,Z_7)=(18623,-2696,-2677,194,393,17,1,8),\qquad Z_k=0\ (k\geq8).
\]
We define symmetric integers $A_{ij}=A_{ji}$ for $i+j\leq N$.  The pure
coefficients are
\[
 A_{0k}=A_{k0}=\operatorname{rnd}(QZ_k/100000).
\]
Let $c_{00}=8$, $c_{10}=c_{01}=2$, $c_{20}=c_{02}=-1$, and $c_{11}=-2$,
with every other $c_{ij}$ zero; after computing $A_{ij}$ set
$C_{ij}=72A_{ij}+Qc_{ij}$.  For $1\leq i\leq j$, in increasing total degree,
form the following sum.  For $0\leq x<i$, $0\leq y<j$, $(x,y)\ne(0,0)$,
put
\[
 h=1+y-x\pmod 3,\quad h\in\{0,1,2\},\qquad
 a=i-x-1+h,\quad c=j-y+1-h,
\]
and set
\[
 w_0(a,c)=c(c-1),\qquad w_1(a,c)=ac,\qquad w_2(a,c)=a(a-1).
\]
Then define
\begin{equation}\label{cert:chart-rec}
 S_{ij}=\sum_{\substack{0\leq x<i,\ 0\leq y<j\\(x,y)\ne(0,0)}}
             C_{xy}A_{a,c}w_h(a,c),\qquad
 A_{ij}=\operatorname{rnd}\left(-\frac{S_{ij}}{C_{00}ij}\right),
 \qquad A_{ji}=A_{ij}.
\end{equation}
The indices $a,c$ are nonnegative, and both $a+c$ and $x+y$ are smaller
than $i+j$.  Thus \eqref{cert:chart-rec} is a definition by finite recursion.
For both choices used here, $21<C_{00}/Q<22$, so its denominator is positive.
Finally put
\begin{equation}\label{cert:P}
 P_{N,b}(z)=2^{-b}\sum_{i+j\leq N}A_{ij}z^i\bar z^j,
 \qquad P=P_{780,500},\qquad P^{\rm ref}=P_{420,224}.
\end{equation}
The chart used in the proof is defined by $P$.  The smaller polynomial
$P^{\rm ref}$ makes the finite evaluations less expensive; uniform bounds
on its difference from $P$ will transfer those evaluations to the chart.
The symmetry of the coefficients makes these polynomials real-valued and
invariant under conjugation.  Rotation permutes the three functions $p_i$;
it need not fix $P$ itself.

\subsection{Uniform residual and polynomial comparisons}
\label{cert:residual}

The reason for the recursion is most transparent at the coefficient level.
For the chart $P_{N,b}$ form
\[
 \mathcal C(z,\bar z)=\sum_{i,j}(C_{ij}/Q)z^i\bar z^j
       =36\bigl[r_0(1-r_0)+2P_{N,b}\bigr].
\]
Write $\mathcal C^{[s]}$ for the terms whose exponent difference is $s$ modulo three.
Then
\[
 \mathcal L=\mathcal C^{[0]}\partial_z\partial_{\bar z}
                 +\mathcal C^{[2]}\partial_z^2+\mathcal C^{[1]}\partial_{\bar z}^2.
\]
The coefficient of $z^{i-1}\bar z^{j-1}$ in $\mathcal L P_{N,b}$ is
\[
 \frac{C_{00}ijA_{ij}+S_{ij}}{Q^2}.
\]
Indeed, the three second derivatives give exactly the shifts and factors in
\eqref{cert:chart-rec}.  A term outside $x<i$, $y<j$ has zero derivative
factor.  Every other coefficient on the right has already been computed.
Nearest rounding therefore leaves a defect of absolute value at most
$C_{00}ij/(2Q^2)<36ij/Q$.  Reflection handles $i>j$.

For any polynomial $F=\sum f_{ij}z^i\bar z^j$, the elementary inequality
\begin{equation}\label{cert:absolute-majorant}
 \sup_{|z|\leq R}|F(z)|\leq\sum_{i,j}|f_{ij}|R^{i+j}
\end{equation}
provides a uniform, rather than sampled, bound.  A convenient majorant for
all low-degree rounding defects is
\[
 \rho=\sum_{k=2}^N\frac{36(k+1)k^2R^{k-2}}Q.
\]
For the remaining products put, for $s=0,1,2$,
\[
 D_{s,k}=\frac{R^k}{Q}\sum_{\substack{a+c=k\\a-c\equiv s\ (3)}}|C_{ac}|,
 \qquad
 V_{s,m}=\frac{R^{m-2}}Q\sum_{a+c=m}|A_{ac}|v_s(a,c),
\]
where $(v_0,v_1,v_2)=(ac,c(c-1),a(a-1))$.  A product with covariance
degree $k$ and differentiated chart degree $m$ has degree $k+m-2$.
Every coefficient outside the recursion is consequently included in
\begin{equation}\label{cert:residual-sum}
 \sup_{|z|\leq R}|\mathcal L P_{N,b}(z)|
 \leq\rho+\sum_{s=0}^2\sum_{k+m>N}D_{s,k}V_{s,m}.
\end{equation}
All terms on the right are nonnegative rational numbers.  They can be
bounded using integer arithmetic: set $U=2^{160}$, replace each $D_{s,k}$
and $V_{s,m}$ by its upper integer multiple at scale $U$, round each
summand of $\rho$ upwards at the same scale, and divide the two sums by
$U^2$ and $U$, respectively.  For $N=780,b=500$ this gives
\begin{equation}\label{cert:epsilon}
 \epsilon:=\sup_{|z|\leq R}|\mathcal L P(z)|
      <3.328667\cdot10^{-22}<3.329\cdot10^{-22}.
\end{equation}
The exact upper numerator and denominator, as well as the complete
coefficients regenerated by \eqref{cert:chart-rec}, form part of the finite
certificate.  Thus the small low-degree defect is attached to the defined
polynomial, not assumed for an arbitrary supplied coefficient list.

For two polynomials, with missing coefficients interpreted as zero, define
\[
 \Delta_a=\sum_{i,j}|P_{ij}-P'_{ij}|R^{i+j}(i+j)^a,
       \qquad a=0,1,2,
\]
using $(i+j)^0=1$ also at the constant term.  A unit real directional
derivative of a monomial costs at most $i+j$ and a second one at most
$(i+j)(i+j-1)$.  Hence
\begin{equation}\label{cert:jet-comparison}
 |P-P'|\leq\Delta_0,\quad
 \|\nabla(P-P')\|\leq\Delta_1/R,\quad
 \|D^2(P-P')\|_{\rm op}\leq\Delta_2/R^2.
\end{equation}
Computing the coefficient differences exactly before taking absolute values,
and rounding the resulting positive sums upwards, gives the following
bounds on the entire closed disk:
\begin{center}
\begin{tabular}{lr}\toprule
Quantity&Upper bound\\\midrule
$|P-P^{\rm ref}|$&$3.517\cdot10^{-17}$\\
$\|\nabla(P-P^{\rm ref})\|$&$1.233\cdot10^{-14}$\\
$\|D^2(P-P^{\rm ref})\|_{\rm op}$&$4.327\cdot10^{-12}$\\
$|P-P^{\rm ref}_{\leq236}|$&$1.902\cdot10^{-11}$\\
$|P-P^{\rm ref}_{\leq220}|$&$7\cdot10^{-11}$\\\bottomrule
\end{tabular}
\end{center}
Here the subscript denotes truncation by total degree.  The same sums with
factors $k$ and $k(k-1)$ prove, for either full or reference chart,
\begin{equation}\label{cert:chart-jets}
 \|\nabla p_i\|<.592,\qquad \|D^2p_i\|_{\rm op}<3.05.
\end{equation}
Truncating either chart at degree 40 changes its value by less than $.0002$
and its gradient by less than $.006$.  These derivative estimates will
control interpolation and root locations; value closeness alone would not
suffice.

\subsection{Geometry of the domain}
\label{cert:geometry}

Work first in the sector $0\leq\theta\leq\pi/3$, writing $z=re^{i\theta}$.
Define its probability domain by $r\leq R$ and $E_1\geq0$.  We shall prove
that the other five coordinates are uniformly positive there.  Rotations
and conjugation then provide the full domain $\Omega$.

From \eqref{cert:chart-jets},
\begin{equation}\label{cert:E-jets}
 \|\nabla E_i\|<G:=\tfrac13+3(.592)<2.11,\qquad
 \|D^2E_i\|_{\rm op}<9.15.
\end{equation}
Divide the sector into $480$ radial and $480$ angular cells, with centers
\[
 r_j=\frac{R(2j+1)}{960},\qquad
 \theta_l=\frac{(2l+1)\pi}{2880},\qquad 0\leq j,l<480.
\]
Evaluate the degree-40 reference chart at these centers, including its
radial derivative and unit-tangential derivative $r^{-1}\partial_\theta$.
The following are the exact integer center tests, after division by the
arithmetic scale:
\begin{align}
 p_i&\geq.0516\quad(i=0,1,2),& E_0,E_2&\geq.055,& d&\geq.488,
       \label{cert:geometry-tests}\\
 E_1<.018&\ \Longrightarrow\ -\partial_rE_1\geq.183
               \quad\hbox{and}\quad\|\nabla E_1\|\geq.3208,\notag\\
 E_1\geq-.008&\ \Longrightarrow\ (9d-2)^2
                 \geq18\sum_{i=0}^2(D_i-D_{i+1})^2.\notag
\end{align}
For completeness, their arithmetic specification uses $T=2^{45}$, the
rational approximation $3141592653589793238/10^{18}$ to $\pi$, and
sine and cosine series through degree 64 on arguments reduced to $[0,2\pi)$.
Every division in these recurrences is rounded down.  The Taylor remainder
is below $10^{-30}$; propagation of rounding errors is bounded by
$65e^{2\pi}/T$.  Including the argument approximation gives error below
$2\cdot10^{-9}$ in either trigonometric value.  There are $861$ possible
chart coefficients through degree 40; their absolute majorant is below
$.921$, and the first-derivative majorant is below $.592$.  Product and
summation errors consequently leave a common allowance $10^{-7}$ for
the displayed scalar fields and gradient vectors.  These rules specify
finite integer evaluations at all $230400$ centers, and all tests
\eqref{cert:geometry-tests} hold.

Let $b=\pi/2880$.  Every point of a cell is within distance
\[
 \delta\leq\sqrt{(R/960)^2+R^2b^2}<.001870298
\]
of its center.  Adding the truncation tails, arithmetic allowance, and
Lipschitz variation gives
\begin{align}
 p_i&>.0516-.0002-.592\delta-10^{-7}>.0502926,
       \label{cert:positive-five}\\
 E_0,E_2&>.055-.0006-G\delta-10^{-7}>.0504548.\notag
\end{align}
If the true $E_1$ at a point is at most $.01$, its truncated center value is
below $.01+.0006+G\delta+10^{-7}<.018$.  Thus the derivative tests apply.
Using the Hessian variation and the degree-40 gradient tail yields
\begin{align}
 -\partial_rE_1&>.183-.018-9.15\delta-Gb-10^{-7}>.1455857,
                 \label{cert:radial-slope}\\
 \|\nabla E_1\|&>.3208-.018-9.15\delta-10^{-7}>.2856866.\notag
\end{align}
Only the first line needs the term $Gb$, accounting for rotation of the
radial unit vector between the point and the center.

In particular $-\partial_rE_1>.144$ wherever $E_1\leq.01$.  At the origin,
$E_1=1/3-P(0)>.1471$.  Once a ray reaches $.01$, its value decreases
strictly while below that level and cannot cross upwards through it.
Consequently each ray has at most one zero.  Let $r_*(\theta)$ be this zero
if it lies below $R$, and $R$ otherwise.  Implicit differentiation on the
smooth face and clipping at $R$ give
\begin{equation}\label{cert:boundary-coarse}
 |r_*'|\leq GR/.144<18.2,\qquad
 |r_*''|\leq\frac{9.15((r_*')^2+R^2)+2.11(2|r_*'|+R)}{.144}<22000,
\end{equation}
where the second inequality is used only on intervals wholly in a smooth
face.  The clipped radius is Lipschitz; its second derivative is not used
across a face--cap junction.

The eigenvalues of the matrix in \eqref{cert:operator} are
\[
 9d\ \pm\ 3\sqrt{2\sum_i(D_i-D_{i+1})^2}.
\]
The center tests give a smaller eigenvalue at least 2; $d\geq.488$ ensures
the required sign before taking square roots.  A cell containing a point
with $E_1\geq0$ has computed center $E_1>-.004546$, so the test is active.
Throughout the sector $r_i\in[-.08,.746667]$, and hence
\[
 \|\nabla D_i\|\leq1.16/3+2(.592)<1.571.
\]
The change from the truncated center to the full value is less than
$2(.0002)+1.571(.001871)+10^{-7}<.00334$.  If all $D_i$ change by at most
$t$, then $\|\Delta C\|_{\rm op}\leq15t$: the three coefficients of $D_i$
in a unit-direction quadratic form are $3+6\cos(2a+2\pi i/3)$; they sum to
9, at most one is negative, and a negative coefficient has magnitude at
most 3.  Their absolute sum is at most 15.  We conclude
\begin{equation}\label{cert:ellipticity}
 \lambda_{\min}(C)>2-15(.00334)>1.949\qquad\hbox{on }\overline\Omega.
\end{equation}
All full/reference differences in \eqref{cert:jet-comparison} fit within
the remaining strict slacks, so these conclusions hold for $P_{780,500}$.

\subsection{The barrier and the interior inequality}
\label{cert:interior}

Finite Taylor polynomials, explicit remainder bounds, and subdivision
into certified boxes are standard techniques of validated numerics; see
\cite[Section~2 and Algorithm~1]{makinoberz2003} and
\cite[Section~2]{solovyevhales2013}.  We specify the integer
arithmetic and all error allowances for the present certificate below.

Let
\[
 \mathcal I=\{(i,j):0\leq j\leq i,\ i+j\leq22,\ i-j\equiv0\pmod3\}.
\]
The $52$ exact rational coefficients $a_{ij}$ in
Table~\ref{cert:barrier-coefficients} define
\begin{equation}\label{cert:barrier}
 \varphi(z)=\Re\sum_{(i,j)\in\mathcal I}a_{ij}z^i\bar z^j.
\end{equation}
For $i>j$, the real part gives
$a_{ij}(z^i\bar z^j+z^j\bar z^i)/2$; a diagonal term occurs once.
This polynomial is invariant under $z\mapsto\zeta z$
and conjugation.  Its constant term is $1.098619383270680538$.

\begin{lemma}[Interior certificate]\label{cert:interior-lemma}
On $\overline\Omega$,
\[
 \mathcal L\varphi>-1+9.8\cdot10^{-8}.
\]
\end{lemma}
\begin{proof}
We specify a finite interval verification and its error accounting.  Let
$q=P^{\rm ref}_{\leq236}$ and $w=z/R$.  By
\eqref{cert:jet-comparison}, $|P-q|<1.902\cdot10^{-11}$ on the disk;
we reserve the looser allowance $10^{-10}$ below.  Define
\[
 A(w)=36\varphi_{z\bar z}(Rw),\qquad B(w)=36\varphi_{zz}(Rw).
\]
Differentiating the real-part expression \eqref{cert:barrier} assigns a
coefficient factor $18R^{i+j-2}$ to each conjugate contribution.  If $i=j$,
the identical contributions combine, correctly counting the original
diagonal term once.

The five real fields to be verified are
\begin{equation}\label{cert:five-fields}
 \mathcal L_q\varphi,\qquad E_{1,q},\qquad A,\qquad \Re B,\qquad\Im B.
\end{equation}
The first supplies the differential lower bound.  The second identifies
boxes lying outside $\Omega$, and the last three bound the error in the
operator when $q$ is replaced by $P$.
The subscript $q$ means that the chart probabilities and covariance are
formed from $q$.  Explicitly, put $r=(1+\Re z)/3$,
$D=r(1-r)+2q$, and let $\Pi_s$ retain exponent differences congruent to
$s$ modulo three.  Then
\[
 d_q=\Pi_0D,\qquad f_q=\Pi_2D,\qquad
 \mathcal L_q\varphi=d_qA+2\Re(f_qB),\qquad
 E_{1,q}=(r+2q)(\zeta z)-3\Pi_0q(z).
\]
All derivative polynomials are evaluated at $w=z/R$.

Convert rational coefficients to integers at scale $U=2^{115}$ by
truncation toward zero.  Divisions in polynomial products are rounded down.
Every constituent field has degree below 320 and fewer than $52000$
possible monomials; the checked coefficient absolute sums, divided by $U$,
are below $2^{36}$.  Before quantization, $2^{36}+1$ is a valid bound.
The sums of coefficient-quantization errors and product-division errors
are therefore below $6\cdot10^{-19}$.  These errors concern the functions
on $|w|\leq1$, by \eqref{cert:absolute-majorant}.

Put $N=128$, $J=14$, $h=\pi/(6N)$.  The boxes
\begin{equation}\label{cert:polar-box}
 z=R\frac{2u+1+v}{2N}\exp\bigl(i(2s+1+t)h\bigr),\qquad
 |t|,|v|\leq1,\quad 0\leq s,u<N,
\end{equation}
cover the whole closed sector.  For each field, expand in $t$ and $v$ to
degree $J$ in each variable.  At scale $T=2^{69}$ the resulting integer
array $a=(a_{ij})_{0\leq i,j\leq14}$ represents
$T^{-1}\sum a_{ij}t^iv^j$.

Here are explicit arithmetic rules for generating these arrays.  Evaluate
$\arctan(1/m)$ by the first 36 terms of its alternating power series,
flooring each unsigned term magnitude at scale $U$ before applying its
alternating sign, and form
\[
 h_U=\left\lfloor\frac{16\arctan_U(1/5)-4\arctan_U(1/239)}{6N}\right\rfloor.
\]
For a monomial $w^i\bar w^j$, let $m=i-j$.  Reduce the phase index
$m(2s+1)$ modulo $12N$, multiply by $h_U$, and evaluate sine and cosine
with terms of orders $0,\ldots,83$.  Generate the order-$k$ angular Taylor
coefficient by differentiating and multiplying by $mh_U/(kU)$, rounding
each division down.  An extra phase index $4Nm$ gives the rotated
constituent in $E_{1,q}$.  Sum the monomial contributions by total degree
and convert once to scale $T$ with floor division.  Finally substitute
$(2u+1+v)/(2N)$ by polynomial Horner evaluation, discarding powers above
$J$ in $v$ and rounding each division down.  These rules produce precisely
the arrays for \eqref{cert:five-fields}.

We next bound what these finite Taylor arrays omit.  For each constituent with integer coefficients $v_{ij}$ at scale $U$
(so its polynomial coefficients are $v_{ij}/U$), the following rational
inequality is checked:
\begin{equation}\label{cert:taylor-tail}
 \sum_{i,j}\frac{|v_{ij}|}{U}\left[
  \frac{((i+j)\pi_+/(6N))^{15}}{15!}
       +4\binom{i+j}{15}(2N)^{-15}\right]<10^{-10},
 \qquad \pi_+=22/7.
\end{equation}
The first term is a real sine/cosine Taylor remainder.  The second bounds
the radial remainder; the angular Taylor polynomial has coefficient sum
at most $\exp(319\pi/(6N))<4$.  Taylor's theorem for the radial monomial
on $[0,1]$ supplies its binomial factor.  The edge field has two
constituents, so its tail bound is doubled.

For clarity, the remaining arithmetic estimates are collected here:
\begin{center}
\begin{tabular}{lr}\toprule
Source of error&Uniform allowance\\\midrule
Coefficient assembly&$6\cdot10^{-19}$\\
Angular recurrence and conversion to scale $T$&$4\cdot10^{-16}$\\
Taylor truncation, edge/other fields&$2\cdot10^{-10}$ / $10^{-10}$\\
Radial Horner floor divisions&$320\cdot15^2/T$\\
At most 28 subdivisions&$28\cdot15^3/T$\\\midrule
Total, for either sign of any field&$3\cdot10^{-10}$\\\bottomrule
\end{tabular}
\end{center}
To justify the angular line, the Machin computation errs by less than
$2/U$ in $h$, and reduced phases by less than $3100/U$.  Taylor recurrence
errors in sine and cosine are below $50000/U$, and each of the 15
angular derivative coefficients errs by less than $300000/U$.  Multiplying
these bounds by the checked coefficient absolute sums gives the stated
conversion allowance.  The radial substitutions contract the coefficient
absolute norm: their affine coefficients are nonnegative and have sum at
most one.  A child substitution has the same contraction property after
absolute values.  Consequently the floor errors add over Horner steps and
subdivision levels rather than amplifying.  There are at most $15^2$
entries per Horner step and fewer than $15^3$ elementary contributions
per subdivision, giving the last two lines.

For an integer array define
\[
 \ell_\sigma(a)=\sigma a_{00}-\sum_{(i,j)\ne(0,0)}|a_{ij}|,
       \qquad\sigma\in\{-1,1\}.
\]
The represented polynomial is bounded below by
$\ell_\sigma(a)/T$ after multiplication by $\sigma$.  Therefore the exact
field satisfies
\begin{equation}\label{cert:jet-lower}
 \sigma F\geq\ell_\sigma(a_F)/T-3\cdot10^{-10}
\end{equation}
everywhere in the box.  If a bound is insufficient, bisect by
$(t,v)=(y,(\varepsilon+x)/2)$, $\varepsilon=\pm1$, thereby alternating
the two subdivision directions.  The child coefficient of $y^ix^k$ is
\[
 \sum_{j\geq k}\left\lfloor
       a_{ij}\binom jk\varepsilon^{j-k}/2^j\right\rfloor.
\]
A box is discarded only when
\begin{equation}\label{cert:outside-test}
 \ell_{-1}(a_{E_{1,q}})>\lfloor T/10^8\rfloor+1.
\end{equation}
Its field error and the chart-change allowance $5\cdot10^{-10}$ total
less than $10^{-8}$, so every such box has full-chart $E_1<0$.

To certify the interior inequality, observe that
\begin{equation}\label{cert:chart-L-error}
 |\mathcal L_P\varphi-\mathcal L_q\varphi|
 \leq2\cdot10^{-10}|A|
              +4\cdot10^{-10}(|\Re B|+|\Im B|).
\end{equation}
For each derivative field $F\in\{A,\Re B,\Im B\}$ and either sign put
$x_{F,\sigma}=-\ell_\sigma(a_F)$.  Use the integer allowance
\[
 e=2\sum_{F\in\{A,\Re B,\Im B\}}
           \sum_{\substack{\sigma=\pm1\\x_{F,\sigma}>0}}
             \left(\left\lfloor\frac{2x_{F,\sigma}}{10^{10}}\right\rfloor+4\right).
\]
The sum of the two positive negative-part bounds dominates the absolute
value of a represented field.  Hence $e/T$ bounds the right-hand side of
\eqref{cert:chart-L-error}, up to less than $2\cdot10^{-18}$ from the
field-approximation errors.  The coefficient used for $A$ is deliberately
larger than necessary.  Each integer cushion exceeds the floor loss in its
summand.

The acceptance condition is
\begin{equation}\label{cert:accept-test}
 \ell_1(a_{\mathcal L_q\varphi})>-T+\lfloor T/10^7\rfloor+e.
\end{equation}
Together with \eqref{cert:jet-lower} and \eqref{cert:chart-L-error}, reserving
$2\cdot10^{-9}$ for all evaluation, chart-change, and threshold-rounding
errors, it proves
$\mathcal L_P\varphi>-1+9.8\cdot10^{-8}$ on the box.

Start with every pair $(s,u)$ in \eqref{cert:polar-box}, apply the two
certification conditions, and recursively subdivide any unresolved box.
All $128$ angular executions terminate with every descendant accepted or
discarded by these conditions at depth at most 28.  The finite certificate
records the complete sector inputs and successful termination for all
sectors, not a collection of floating-point sample values.  Arithmetic can
be interpreted as exact integers with the specified floor divisions.
For the fixed-width implementation, array norms are below $2^{109}$,
Horner products below $2^{117}$, and subdivision products below $2^{121}$,
within signed 128-bit range.  Section~\ref{cert:reproducibility} records the
machine-semantics checks and reproduced finite coverage.  Rotations and
conjugation cover the other sectors.  This proves the lemma.
\end{proof}

\subsection{The six-component boundary inequality}
\label{cert:boundary}

We apply the three-state seed to a face of the six-component domain, and
the matching construction to its circular cap.  It is important to keep
these two sources of attainable values distinct.

Put $U=U_{\gamma_0}$, where $\gamma_0=300000/225925$.  Since
$3/\beta>\gamma_0$, Proposition~\ref{prop:monotonicity} permits us to use
this lower parameter in the face estimate of Proposition~\ref{prop:face}.
At $E_1=0$ define
\begin{equation}\label{cert:face-data}
 Q=(p_2+E_0,p_0+E_2,p_1),\qquad
 J=\en(p_2)+\en(p_0)+\en(E_0)+\en(E_2)-\en(Q_0)-\en(Q_1).
\end{equation}
The face lower bound is
\begin{equation}\label{cert:Be}
 B_e=\frac{1129}{1500}\bigl(U(Q)+J\bigr).
\end{equation}
At the cap put $l_h=E_{h+1}+E_{h+2}$.  The side-$h$ marginal is
$(1-l_h-p_h,l_h,p_h)$, so Proposition~\ref{prop:matching} gives
\begin{equation}\label{cert:Bc}
 B_c=\frac13\sum_{h=0}^2
       \bigl[\en(l_h)+\en(p_h)+\en(1-l_h-p_h)\bigr].
\end{equation}

\paragraph{A finite lower estimate for $U$.}
For evaluation in \eqref{cert:Be}, use the entropy baseline, the improved
shell bound, and finite-library mixtures from
Section~\ref{cert:w-seed-section}.
The library bounds are tabulated on the three-state simplex grid of spacing
$1/500$.  Within a grid square whose four vertices are in the simplex,
bilinear interpolation is a convex combination of the four vertex laws
and reproduces the target law exactly.  Concavity of $U$ therefore makes
the interpolant of lower utilities a valid lower bound.  Every grid vertex
used here has all coordinates greater than $.03$.

At each grid vertex subtract $2\cdot10^{-6}$ from the computed library
utility.  Subtract a further $2\cdot10^{-8}$ from the maximum of this
interpolant and $F_0$.  Inside the shell also consider
$F_0+H_W-24\cdot10^{-6}$, using a radial guard $10^{-8}$ to ensure that the
point really lies in the shell.  Take the maximum of the available bounds
and cap it at $1.45$.  The resulting evaluated estimate is nonnegative.
The error recurrences, entropy evaluation, and distribution correction
proved above allow at most $10^{-7}$ additional overstatement; we include
it in the boundary arithmetic allowance.  Neither selection of the maximum
nor downward capping requires the library to be optimal.

\paragraph{Locating boundary points.}
Set
\[
 K=65536,\qquad h=\frac{\pi}{3K}.
\]
Evaluate the boundary at the $K+1$ endpoint angles.  Use the degree-220
reference chart and arithmetic scale $B=2^{52}$.  The value error relative
to $P$ is below $7\cdot10^{-11}$, and each $E$-coordinate error is below
$2.5\cdot10^{-10}$.  The chart coefficient tail was already bounded in
Section~\ref{cert:residual}.  For evaluation, combine conjugate coefficients
as real cosines and evaluate radial powers by Horner's rule.  The fundamental cosine is computed by the degree-48 Taylor polynomial
on $[0,\pi]$ after symmetry reduction, with error at most $300/B$.
Higher angular harmonics are formed by the recurrence
$c_{s+1}=2c_1c_s-c_{s-1}$; their error amplification is controlled by
the $k^2$-weighted coefficient sum.  The total absolute coefficient
sum of the reference chart on the disk is below $.414$, and its
$k^2$-weighted sum is below 6.  These majorants bound the additional
coefficient, trigonometric, and Horner errors within the stated
allowances.

If the computed $E_1(Re^{i\theta})$ is nonnegative, take the clipped
radius $R$; otherwise perform 38 bisections on $[0,R]$.  By
\eqref{cert:radial-slope}, the combined evaluation and bisection error in
$r_*(\theta)$ is below $2\cdot10^{-9}$: the dominant term is
$(2.5\cdot10^{-10})/.144$, and $R2^{-38}$ is smaller than the remaining
slack.  The same bound holds when a sign uncertainty changes whether the
root is clipped.

Use the band
\begin{equation}\label{cert:band}
 b_K=2.7h+5\cdot10^{-9}
\end{equation}
to classify a complete angular cell from its left endpoint.  This is
valid because $|\partial_\theta E_1(Re^{i\theta})|<GR<2.62$.  A value
below $-b_K$ identifies a full face cell; one above $b_K$ identifies a
full cap cell.  In the other cells, subdivide into $8192$ equal subcells
and use the analogous band at that scale.  Where the sign remains
uncertain, evaluate both bounds.

At an endpoint the face bound is computed whenever the cap value is
below $2.1b_K$, and the cap bound whenever it is above $-2.1b_K$.
Such a face evaluation may occur at a clipped cap point with $E_1\ne0$.
Its three-state law is normalized by setting
\[
 \widehat Q_0=p_2+E_0,\qquad
 \widehat Q_1=p_0+E_2,\qquad
 \widehat Q_2=1-\widehat Q_0-\widehat Q_1=p_1+E_1.
\]
On the face, $\widehat Q=Q$.  At a sampled point off the face, use
$\widehat Q$ in the right side of \eqref{cert:Be}, with the same
internal coordinates and first two coarse coordinates in
\eqref{cert:face-data}.  This is an auxiliary comparison expression;
the interpolation and displacement estimates below transport it to actual
face laws, where Proposition~\ref{prop:face} supplies attainability.

The band rule ensures that both endpoints required for a full-cell test have been
evaluated: for example, a left value below $-b_K$ forces the next value
below $-b_K+2.62h$ plus evaluation error, which is below $2.1b_K$.
This also guarantees the face samples near cap points at which the inset
will become active.  There are six transition cells in the finite check.

\paragraph{Barrier evaluation and derivative bounds.}
For \eqref{cert:barrier}, store the scaled radial coefficients at
$B_\varphi=2^{44}$, rounded to nearest integers.  Evaluate powers of
$r/R$ and the appropriate cosine or sine for angular differentiation.
For a term of radial degree $d=i+j$ and angular frequency $s=i-j$, an
$a$-th radial and $b$-th angular derivative has multiplier
$d(d-1)\cdots(d-a+1)s^b$ with the appropriate phase shift.  Exact
coefficient sums give the following value and derivative error bounds:
\begin{equation}\label{cert:barrier-arithmetic}
 \text{value error}<1.2\cdot10^{-7},\qquad
 \text{error in each derivative through order three}<.002.
\end{equation}
One may obtain these by summing, for each monomial, its coefficient
rounding error, the at most $d$ radial-power division errors, and its
trigonometric error, multiplied by the displayed derivative factor.
The absolute radial coefficient sums weighted by $d^k$, $k=0,\ldots,4$,
are respectively below
\[
 35000,\quad750000,\quad15000000,\quad300000000,\quad6000000000.
\]
These conservative rational bounds make the calculation independent of
cancellation in the evaluated polynomial.

At every sampled point where a face bound is evaluated, the integer
checks in physical polar coordinates give
\begin{equation}\label{cert:face-jet-tests}
 |\partial_r^a\partial_\theta^b\varphi|<
 \begin{cases}2.9,&a+b=1,\\79,&a+b=2,\\1999,&a+b=3.\end{cases}
\end{equation}
When evaluated using normalized radius $r/R$, the threshold contains
$R^a$, converting it to precisely this physical-coordinate statement.
Cap samples satisfy $|\varphi_\theta|<999$ and
$|\varphi_{\theta\theta}|<4999$ before the errors
\eqref{cert:barrier-arithmetic} are restored.

A useful uniform fourth-derivative bound is
\begin{equation}\label{cert:fourth}
 \sup_{r\leq R}|\partial_r^a\partial_\theta^b\varphi|
 \leq\sum_{\substack{(i,j)\in\mathcal I\\i+j\geq a}}
 |a_{ij}|\frac{(i+j)!}{(i+j-a)!}R^{i+j-a}|i-j|^b
 <435\cdot10^6\qquad(a+b=4).
\end{equation}
These are five finite rational sums.  Their largest value is less than
$434439204$, so the stated integer bound has spare margin.  Taylor's
theorem with \eqref{cert:fourth} extends the sampled derivative estimates
between boundary samples.

\paragraph{Interpolation on a smooth face.}
On a full face cell the additional integer secant test is
\[
 \frac{|\widehat r_*(\theta+h)-\widehat r_*(\theta)|}{h}<3.05,
\]
where hats denote the computed physical radii.  The coarse bound
$|r_*''|<22000$ and the root-location errors show
$|r_*'|<3.05+22000h+4\cdot10^{-9}/h<3.41$.  Substitute this in
\eqref{cert:boundary-coarse} to get
\begin{equation}\label{cert:boundary-fine}
 |r_*'|<3.41,\qquad |r_*''|<1000.
\end{equation}
Write $z(\theta)=r_*(\theta)e^{i\theta}$.  Then
\[
 |z'|^2\leq3.41^2+R^2,\qquad |z''|\leq1000+2(3.41)+R.
\]
Equations~\eqref{cert:chart-jets} and \eqref{cert:E-jets} now give
\begin{equation}\label{cert:face-curvature}
 |Q_0''|,|Q_1''|<2900,\qquad |J''|<30000,\qquad
 |(\varphi\circ z)''|<5100.
\end{equation}
Here are details for the last two estimates.  In differentiating $\en(s)$,
use $\en'(s)=-1-\log s$ and $\en''(s)=-1/s$.
The four internal arguments in $J$ exceed $.05$, and $Q_0,Q_1>.10$.
The first-derivative bounds for an internal $p$ and $E$ are
$.592|z'|$ and $2.11|z'|$; their second derivatives are bounded by
$3.05|z'|^2+.592|z''|$ and $9.15|z'|^2+2.11|z''|$.
The six entropy chain-rule terms are bounded using
$|\en'(s)|<2$ on $[.05,1]$ and $|\en'(s)|<1.31$ on $[.10,1]$.
Their sum is below 30000.  For the barrier, expand its first and second
polar derivatives about a sampled endpoint, using
\eqref{cert:face-jet-tests}, \eqref{cert:barrier-arithmetic}, and
\eqref{cert:fourth}.  The polar displacement is at most
$(3.41+1)h+4\cdot10^{-9}$.  Throughout the cell the first derivatives
are below $2.91$ and the second derivatives below $80.3$.  Therefore
\[
 |(\varphi\circ z)''|
 \leq80.3(3.41+1)^2+2.91(1000)<5100.
\]

For a function with $|f''|\leq M$, its difference from linear
interpolation over an interval of length $h$ is at most $Mh^2/8$.
Concavity of $U$ lets us interpolate the closed attainable law--utility
pairs at neighboring endpoints, giving a lower bound at the chord law.
The first two $Q$ coordinates differ from the curved target by at most
$2900h^2/8$.  On the actual face the third coordinate is $Q_2=p_1$;
using it directly gives
\[
 |Q_2''|\leq3.05(3.41+1.24)^2
            +.592(1000+2(3.41)+1.24)<663.
\]
Since $Q_0,Q_1>.10$ and $Q_2>.05$, the curvature divided by a coordinate
lower bound is at most $29000$.  The exact-law correction of
Lemma~\ref{cert:w-correction} costs at most $1.5\cdot29000h^2/8$ in $U$.
Adding the conditional-entropy
and barrier interpolation costs gives
\begin{equation}\label{cert:face-interpolation}
 \left[5100+.754\bigl(30000+1.5(29000)\bigr)\right]\frac{h^2}{8}
       <2\cdot10^{-6},
\end{equation}
where $.754>1129/1500$.

\paragraph{Interpolation on a cap.}
On an actual cap all entropy arguments in \eqref{cert:Bc} exceed $.049$.
The chart derivative bounds and $|z'|=|z''|=R$ give
$|B_c''|<1600$.  For example, the gradient bounds for the three marginal
arguments are $4.22$, $.592$, and $4.812$, and the Hessian bounds are
$18.3$, $3.05$, and $21.35$; substituting these in the entropy chain rule
already gives the stated bound.  The cap sample derivative checks extend
to $|\varphi_{\theta\theta}|<10000$.  A direct coefficient bound gives
$|\varphi_{\theta\theta\theta}|<3875770$, so the second derivative
can move by less than 62 in one full cell.  It follows that the full-cap
interpolation loss is less than
$(1600+10000)h^2/8<.4\cdot10^{-6}$.

\paragraph{Transition subcells.}
No second derivative across the clipped corner is needed.  In a subcell
the actual boundary point and its sampled clipped point have complex
displacement less than
\[
 (18.2+R)h/8192+2\cdot10^{-9}<4.1\cdot10^{-8}.
\]
This changes $Q_0,Q_1$ by less than $1.12\cdot10^{-7}$ each, and their
reconstructed third coordinate by less than $2.24\cdot10^{-7}$.
The singleton correction costs at most
\[
 \frac{1129}{1500}\,1.5\,
       \max\left\{\frac{1.12\cdot10^{-7}}{.10},
                    \frac{2.24\cdot10^{-7}}{.05}\right\}
       <5.06\cdot10^{-6}.
\]
The conditional-entropy differential for one split is
$\log((a+b)/a)\,da+\log((a+b)/b)\,db$.  All internal coordinates
remain above $.049$ on this displacement, so each logarithmic factor
is less than $3.1$.  Summing the four internal-coordinate changes
and multiplying by $1129/1500$ gives less than $.52\cdot10^{-6}$.
The local barrier first derivatives, extended from
\eqref{cert:face-jet-tests} with \eqref{cert:fourth}, are below 3;
its change is less than $.13\cdot10^{-6}$.  Thus the total face
\emph{displacement} loss is less than $6.5\cdot10^{-6}$.
This estimate remains valid if the sampled clipped point is just on the
cap: reconstructing $Q_2$ absorbs its small active-coordinate mass into
the five-state face law, and the same correction reaches the actual
face law.

For a cap transition, evaluate the smooth entropy expression in
\eqref{cert:Bc} even if the sampled cap point has $E_1<0$.  Attainability
is asserted only at the actual cap target, where Proposition~\ref{prop:matching}
applies; the displacement estimate compares the two entropy expressions.
All their arguments stay above $.049$, and $|B_c'|<26$.  The cap derivative checks, followed by the local second
bound just established, give $|\varphi_\theta|<1001$.  The subcell
width is below $2\cdot10^{-9}$, so displacement together with the cap
endpoint arithmetic allowance costs less than $2.4\cdot10^{-6}$.

\paragraph{Arithmetic checks and retained margins.}
The finite evaluations compare the face and cap lower estimates against
$\varphi$.  The smallest recorded integer gaps at scale $B=2^{52}$ are
\begin{equation}\label{cert:actual-gaps}
 g_e=128315602072/B>28.49\cdot10^{-6},\qquad
 g_c=152568202355/B>33.87\cdot10^{-6}.
\end{equation}
All $65536$ cells, and all $8192$ subcells of each of the six transition
cells, satisfy the corresponding tests.  On full face cells the secant
tests used above also hold.  The conditional laws and mixture checks
used to produce these values are the finite rational constructions of
Section~\ref{cert:w-seed-section}.

A conservative endpoint arithmetic allowance is $10\cdot10^{-6}$ for
a face evaluation and $.2\cdot10^{-6}$ for a cap evaluation.  To see
that the former suffices, account separately for the $2\cdot10^{-9}$
root-location error with the chart gradient bounds, correct the resulting
coarse-law error by singleton dilution with coordinates above $.03$,
include the $10^{-7}$ possible error in the computed lower estimate for
$U$, and add the entropy and barrier errors.  Each coordinate moves by
less than $6\cdot10^{-9}$, with less than $12\cdot10^{-9}$ for the
reconstructed coordinate; this dilution costs less than $6\cdot10^{-7}$
before the prefactor.  The entropy and barrier evaluation estimates above
fit well within the remaining allowance.  A cap needs no root correction;
the chart error in its entropy arguments and the barrier value error give
the stated $.2\cdot10^{-6}$ bound.  All trigonometric, logarithmic,
coefficient, and integer-division errors have been included.

We use the recorded face minimum in \eqref{cert:actual-gaps} and charge
the full face arithmetic allowance in transition cells separately from
displacement.  The resulting bounds are
\begin{center}
\begin{tabular}{lrrr}\toprule
Boundary case&Gap used&Arithmetic loss&Interpolation/displacement loss\\\midrule
Full face&$28.49\cdot10^{-6}$&$10\cdot10^{-6}$&$2\cdot10^{-6}$\\
Transition face&$28.49\cdot10^{-6}$&$10\cdot10^{-6}$&$6.5\cdot10^{-6}$\\
Full cap&$8\cdot10^{-6}$&$.2\cdot10^{-6}$&$.4\cdot10^{-6}$\\
Transition cap&$8\cdot10^{-6}$&\multicolumn{2}{c}{$2.4\cdot10^{-6}$ combined}\\\bottomrule
\end{tabular}
\end{center}
Every entry leaves more than $5.5\cdot10^{-6}$.

All arithmetic has an exact-integer interpretation with specified
truncating or floor divisions.  In the finite-width implementation,
probabilities and counts are below $3B$, and scores and accumulated
utilities below $100B$.  Projected-law determinants and barycentric
numerators have magnitude at most $2^{60}$; utility-gap products are
below $2^{94}$, ratio-test products below $2^{122}$, and barrier
products below $10^{36}<2^{127}$.  Narrowing casts are checked.
These bounds justify the integer interpretation independently of the
additional runtime checks described in Section~\ref{cert:reproducibility}.
We have proved the following uniform statement.

\begin{lemma}[Continuous boundary margin]\label{cert:boundary-lemma}
On the boundary of the full six-state probability domain,
\[
 F_\beta(\alpha(z))\geq\varphi(z)+5.5\cdot10^{-6}.
\]
\end{lemma}

\subsection{The positive inset and the comparison argument}
\label{cert:inset}

The supporting-covector estimate requires all six probabilities to be
bounded away from zero.  Set $\eta=2\cdot10^{-9}$ and replace each active
face $E_i=0$ by $E_i=\eta$, retaining the cap where it is reached first.
Let the resulting domain be $\Omega_\eta$.  It contains the origin,
has compact closure, and all six coordinates are at least $\eta$ there.
We now transfer Lemma~\ref{cert:boundary-lemma} to its boundary.

On a removed terminal segment of a ray,
$0\leq E_1\leq\eta$ and $-\partial_rE_1>.144$.  Its length is at most
\begin{equation}\label{cert:inset-distance}
 \Delta r\leq\eta/.144<1.389\cdot10^{-8}.
\end{equation}
This also handles an original cap endpoint with $0\leq E_1<\eta$.
The five other coordinates remain greater than $.05$ and each changes
by less than $2.11\Delta r<2.94\cdot10^{-8}$.  Write $\alpha_b$ and
$\alpha_\eta$ for the original and new endpoint laws.  With
\[
 \lambda=\max_{i:(\alpha_b)_i>0}
             \left(1-\frac{(\alpha_\eta)_i}{(\alpha_b)_i}\right)_+,
\]
the active coordinate contributes nothing, and
$\lambda<2.94\cdot10^{-8}/.05=5.88\cdot10^{-7}$.
If $\lambda>0$, the vector
\[
 s=\frac{\alpha_\eta-(1-\lambda)\alpha_b}{\lambda}
\]
is a probability distribution.  Mix a construction at $\alpha_b$ with
singleton constructions of average law $s$.  Singleton utility is
nonnegative, and the original utility is at most $\log6<2$.
The loss is below $2\lambda<1.18\cdot10^{-6}$.  When $\lambda=0$ no
correction is required.  The closure of rational mixtures gives the
same conclusion for limiting endpoint laws.

To bound the change of the barrier, every point of the inset segment has
a face-check sample within angular distance $h=\pi/(3\cdot65536)$.
The band test covers even cap endpoints where the inset becomes active.
The radial Lipschitz bound, root error, and inset displacement give total
polar-coordinate distance at most
\[
 d_*=19.2h+2\cdot10^{-9}+1.389\cdot10^{-8}<.000307.
\]
Taylor's theorem for $\partial_r\varphi$, using
\eqref{cert:face-jet-tests}, \eqref{cert:barrier-arithmetic}, and
\eqref{cert:fourth}, gives
\[
 |\partial_r\varphi|
 <2.902+79.002d_*+\frac{1999.002}{2}d_*^2
                         +\frac{435\cdot10^6}{6}d_*^3<3.
\]
Thus the barrier changes by less than $3\Delta r<4.2\cdot10^{-8}$.
The two transfer losses total less than $1.3\cdot10^{-6}$, leaving
more than $4.2\cdot10^{-6}$ slack on the inset boundary.  At unchanged
cap points there is no loss.  Symmetry treats all faces.

\begin{proof}[Proof of Theorem~\ref{thm:certificate}]
Let $g$ be a supporting covector of $F_\beta$ at a point of the inset,
normalized by $g\cdot\alpha=F_\beta(\alpha)$.  Nonnegativity at simplex
vertices implies $g_i\geq0$.  Proposition~\ref{prop:values} gives
$F_\beta\leq\log6<2$.  Rotation covariance in
\eqref{ap:complex-operator} and the residual bound \eqref{cert:epsilon} imply
$|\mathcal Lp_i|\leq\epsilon$ and
$|\mathcal LE_i|\leq3\epsilon$.  Lemma~\ref{ap:chart-error} consequently
gives
\begin{equation}\label{cert:covector-debt}
 |g\cdot\mathcal L\alpha|
 \leq2\max_i\frac{|\mathcal L\alpha_i|}{\alpha_i}
 \leq\frac{6\epsilon}{\eta}
 <9.987\cdot10^{-13}<10^{-12}.
\end{equation}
The chart and value function are continuous on this compact inset.
If $\varphi-F_\beta\circ\alpha$ had a positive maximum, the boundary
slack would place it at an interior point.  Subtracting the maximum from
$\varphi$ gives a smooth lower test there.  Theorem~\ref{thm:ap},
Lemma~\ref{cert:interior-lemma}, and \eqref{cert:covector-debt} would imply
\[
 1\leq g\cdot\mathcal L\alpha-\mathcal L\varphi
       <10^{-12}+1-9.8\cdot10^{-8}<1,
\]
a contradiction.  Therefore $F_\beta(\alpha(z))\geq\varphi(z)$ throughout
the inset.  At the origin $\alpha(z)=\alpha_*$, and
\[
 \varphi(0)-\log3
 =1.098619383270680538-\log3>7.094\cdot10^{-6}>0.
\]
For example, the logarithm inequality follows from its positive
$2\sum_{j\geq0}t^{2j+1}/(2j+1)$ series with $t=(3-1)/(3+1)=1/2$ and
a geometric bound on the tail; no floating-point comparison is needed.
This proves the theorem.
\end{proof}
